\chapter{Introduction}
\label{ch:introduction}

Card payments make up a large and growing share of everyday spending, and card fraud has
grown with them. Banks and payment processors lose money to fraud every year, and
cardholders lose time and trust when their cards are misused. To limit these losses,
payment systems score transactions as they arrive and flag suspicious ones for review or
blocking. Early systems used fixed rules written by analysts. Rules are easy to read but
slow to keep up with new fraud patterns, so most modern systems now use \gls{ml} models
that learn patterns directly from past transactions.

\section{Background}

A fraud detection model is only useful if it fits the shape of the problem, and this
problem has three properties that make it hard.

First, fraud is extremely rare. In the data used here, fraudulent transactions
make up 0.167 percent of all records, or roughly one in every 600. A model trained without
regard for this imbalance will lean toward calling everything legitimate, because that is
correct 99.83 percent of the time. This also makes accuracy useless as a headline metric,
and makes even the \gls{roc} curve optimistic, because the huge pool of correctly cleared
transactions dominates the false alarm rate
\citep{davis2006relationship,saito2015precision}.

Second, the features are anonymised. To protect privacy, the public dataset
releases most fields as \gls{pca} components rather than raw values, so feature engineering
based on domain knowledge is not possible. We can see which components carry signal, but
not what those components mean.

Third, the cost of a mistake is asymmetric. Missing a fraud is usually worse than
raising a false alarm, and the two errors are measured in different units. A missed fraud
costs the transaction amount, while a false alarm costs an analyst's time. Where the model
draws its decision line is therefore a business decision, not an implementation detail.

\section{Motivation}

Much published work on this dataset reports a single headline number from a single
train/test split. That practice is fragile. With only 95 frauds in a 20 percent test set,
a handful of transactions moving across the decision boundary can shift recall by several
percentage points. A central motivation of this project is to check whether the differences
between competing models survive basic statistical testing, using confidence intervals,
paired significance tests, and cross-validation. On this dataset those differences turn out
to be mostly noise, and reporting them as findings would be misleading.

\section{Objectives}

The primary objectives of this project are:

\begin{enumerate}
  \item To build a leakage-free pipeline and use it to train and compare eight model
        pipelines, being four algorithms each trained with class weighting and with
        \gls{smote}, ranked by \gls{prauc}, the metric suited to a rare positive class.
  \item To measure the uncertainty in those results using bootstrap confidence intervals,
        pairwise McNemar tests, and $k$-fold cross-validation, and to state plainly which
        apparent differences are real.
  \item To serve the trained models through a FastAPI backend and an interactive React web
        application in which a reviewer can score transactions, switch models, and move the
        decision threshold.
\end{enumerate}

\section{Scope and Significance}

The project is built around a single public dataset and works offline: models are trained
once and then served. It does not attempt streaming ingestion or online retraining. Within
that scope the aim is to get the fundamentals right rather than to chase the highest
possible score. The result is a working system that is clear about its limits. It reports
how many frauds it still misses, explains why accuracy is not used, puts an uncertainty
range around every headline figure, and exposes the decision threshold as a live control so
a reviewer can see the trade-off rather than be told about it.

\section{Contributions}

\begin{itemize}
  \item A leakage-free pipeline in which de-duplication happens before splitting and all
        resampling happens inside the training fold.
  \item A comparison of eight pipelines with the uncertainty measured: bootstrap intervals
        on PR-AUC, Wilson intervals on precision and recall, and Holm-corrected McNemar
        tests on every pair.
  \item A controlled resampling study that separates the sampling step from the algorithm.
  \item A calibration, cost, lift, and learning-curve analysis of the deployed model.
  \item A served system with a model registry and an interactive web application, including
        a 3D rendering of the full decision surface across all eight models and 201
        thresholds.
  \item A reproducible reporting chain: every table in this document is generated
        automatically from the analysis outputs, so no number in the text can drift away
        from the code that produced it.
\end{itemize}
